\documentclass{article}
\usepackage[utf8]{inputenc}
\usepackage[italian]{babel}
\usepackage{amsmath, amssymb}
\usepackage{tikz}

\begin{document}

\section*{SOVRAPPOSIZIONE DI MOTI ARMONICI}

A) \textbf{moti armonici in direzioni parallele ($\hat{x}$)}
\begin{equation*}
\begin{cases}
x_1(t) = A_1 \cos(\omega_1 t + \varphi_1) \\
x_2(t) = A_2 \cos(\omega_2 t + \varphi_2)
\end{cases}
\implies x(t) = A_1 \cos(\omega_1 t + \varphi_1) + A_2 \cos(\omega_2 t + \varphi_2)
\end{equation*}

1) $\omega_1 = \omega_2 = \omega$ (vediamo che $\exists (A, \varphi) \mid x(t) = A \cos(\omega t + \varphi)$)
\begin{equation*}
x(t) = A_1 \cos(\omega t + \varphi_1) + A_2 \cos(\omega t + \varphi_2)
\end{equation*}
\begin{align*}
A \cos(\omega t + \varphi) &= A [\cos(\omega t)\cos\varphi - \sin(\omega t)\sin\varphi] \\
&= A \cos\varphi \cos(\omega t) - A \sin\varphi \sin(\omega t)
\end{align*}
\begin{align*}
&A_1 \cos(\omega t + \varphi_1) + A_2 \cos(\omega t + \varphi_2) = \\
&= A_1 [\cos(\omega t)\cos\varphi_1 - \sin(\omega t)\sin\varphi_1] + A_2 [\cos(\omega t)\cos\varphi_2 - \sin(\omega t)\sin\varphi_2] \\
&= (A_1 \cos\varphi_1 + A_2 \cos\varphi_2)\cos(\omega t) - (A_1 \sin\varphi_1 + A_2 \sin\varphi_2)\sin(\omega t)
\end{align*}

\begin{equation*}
\implies A \cos\varphi = A_1 \cos\varphi_1 + A_2 \cos\varphi_2 \quad \text{e} \quad A \sin\varphi = A_1 \sin\varphi_1 + A_2 \sin\varphi_2
\end{equation*}
\begin{equation*}
\implies \text{tg}\,\varphi = \frac{A_1 \sin\varphi_1 + A_2 \sin\varphi_2}{A_1 \cos\varphi_1 + A_2 \cos\varphi_2} \implies A = \sqrt{A_1^2 + A_2^2 + 2A_1 A_2 \cos(\varphi_1 - \varphi_2)}
\end{equation*}

\begin{align*}
\implies (A \cos\varphi)^2 + (A \sin\varphi)^2 = A^2 &= (A_1 \cos\varphi_1 + A_2 \cos\varphi_2)^2 + (A_1 \sin\varphi_1 + A_2 \sin\varphi_2)^2 \\
&= A_1^2 + A_2^2 + 2A_1 A_2 (\cos\varphi_1 \cos\varphi_2 + \sin\varphi_1 \sin\varphi_2) \\
&= A_1^2 + A_2^2 + 2A_1 A_2 \cos(\varphi_1 - \varphi_2)
\end{align*}

\subsection*{1.1) $\varphi_1 = \varphi_2$ CONCORDANZA DI FASE / INTERFERENZA COSTRUTTIVA}
\begin{equation*}
\implies A = A_1 + A_2 \qquad \implies \text{tg}\,\varphi = \text{tg}\,\varphi_1 = \text{tg}\,\varphi_2
\end{equation*}
% Grafico concordanza di fase
\begin{tikzpicture}[scale=0.8]
\draw[->] (-0.2,0) -- (5,0) node[right] {$x$};
\draw[->] (0,-1.5) -- (0,1.5);
\draw[blue, thick, domain=0:4.5, samples=100] plot (\x, {sin(\x*180/pi)});
\draw[green!60!black, thick, domain=0:4.5, samples=100] plot (\x, {0.5*sin(\x*180/pi)});
\draw[red, thick, domain=0:4.5, samples=100] plot (\x, {1.5*sin(\x*180/pi)});
\node at (2.5, 1.8) {\tiny $A=A_1+A_2$};
\node at (2.5, -0.7) {\tiny $A_1$};
\node at (2.5, -1.2) {\tiny $A_2$};
\end{tikzpicture}

\subsection*{1.2) $\varphi_1 = \varphi_2 + \pi$ OPPOSIZIONE DI FASE / INTERFERENZA DISTRUTTIVA}
\begin{equation*}
\implies A = |A_1 - A_2| \qquad \implies \text{tg}\,\varphi = \text{tg}\,\varphi_1 = \text{tg}\,\varphi_2
\end{equation*}
% Grafico opposizione di fase
\begin{tikzpicture}[scale=0.8]
\draw[->] (-0.2,0) -- (5,0) node[right] {$x$};
\draw[->] (0,-1.5) -- (0,1.5);
\draw[blue, thick, domain=0:4.5, samples=100] plot (\x, {sin(\x*180/pi)});
\draw[green!60!black, thick, domain=0:4.5, samples=100] plot (\x, {-0.5*sin(\x*180/pi)});
\draw[red, thick, domain=0:4.5, samples=100] plot (\x, {0.5*sin(\x*180/pi)});
\node at (2.5, 0.8) {\tiny $A=|A_1-A_2|$};
\end{tikzpicture}

\subsection*{1.3) $\varphi_1 = \varphi_2 + \pi/2$ QUADRATURA}
\begin{equation*}
\implies A = \sqrt{A_1^2 + A_2^2} \qquad \implies \text{tg}\,\varphi = \frac{A_1 + \text{tg}\,\varphi_2 A_2}{A_2 - \text{tg}\,\varphi_2 A_1}
\end{equation*}
% Grafico in quadratura
\begin{tikzpicture}[scale=0.8]
\draw[->] (-0.2,0) -- (5,0) node[right] {$x$};
\draw[->] (0,-1.5) -- (0,1.5);
\draw[blue, thick, domain=0:4.5, samples=100] plot (\x, {sin(\x*180/pi)});
\draw[green!60!black, thick, domain=0:4.5, samples=100] plot (\x, {0.8*cos(\x*180/pi)});
\draw[red, thick, domain=0:4.5, samples=100] plot (\x, {sqrt(1+0.64)*sin(\x*180/pi + 38.6)});
\end{tikzpicture}

\subsection*{1.4) $A_1 = A_2$}
\begin{equation*}
\implies A = 2A_1 \left| \cos\left(\frac{\varphi_1 - \varphi_2}{2}\right) \right| \qquad \implies \text{tg}\,\varphi = \frac{\sin\varphi_1 + \sin\varphi_2}{\cos\varphi_1 + \cos\varphi_2}
\end{equation*}

Quando $A_1 = A_2$ possiamo avere i 3 casi precedenti (1.1, 1.2, 1.3):

\paragraph*{1.4.1) $A_1 = A_2$, $\varphi_1 = \varphi_2 \implies A = 2A_1$}
$\hookrightarrow$ IN FASE
% Grafico
\begin{tikzpicture}[scale=0.5]
\draw[->] (-0.2,0) -- (5,0) node[right] {$x$};
\draw[->] (0,-1.5) -- (0,1.5);
\draw[blue, thick, domain=0:4.5, samples=100] plot (\x, {0.6*sin(\x*180/pi)});
\draw[red, thick, domain=0:4.5, samples=100] plot (\x, {1.2*sin(\x*180/pi)});
\end{tikzpicture}

\paragraph*{1.4.2) $A_1 = A_2$, $\varphi_1 = \varphi_2 + \pi \implies A = 0$}
$\hookrightarrow$ IN CONTROFASE
% Grafico
\begin{tikzpicture}[scale=0.5]
\draw[->] (-0.2,0) -- (5,0) node[right] {$x$};
\draw[->] (0,-1.5) -- (0,1.5);
\draw[blue, thick, domain=0:4.5, samples=100] plot (\x, {0.6*sin(\x*180/pi)});
\draw[green!60!black, thick, domain=0:4.5, samples=100] plot (\x, {-0.6*sin(\x*180/pi)});
\draw[red, thick, domain=0:4.5, samples=100] plot (\x, {0});
\end{tikzpicture}

\paragraph*{1.4.3) $A_1 = A_2$, $\varphi_1 = \varphi_2 + \pi/2 \implies A = \sqrt{2}A_1$}
$\hookrightarrow$ IN QUADRATURA
% Grafico
\begin{tikzpicture}[scale=0.5]
\draw[->] (-0.2,0) -- (5,0) node[right] {$x$};
\draw[->] (0,-1.5) -- (0,1.5);
\draw[blue, thick, domain=0:4.5, samples=100] plot (\x, {0.6*sin(\x*180/pi)});
\draw[green!60!black, thick, domain=0:4.5, samples=100] plot (\x, {0.6*cos(\x*180/pi)});
\draw[red, thick, domain=0:4.5, samples=100] plot (\x, {0.6*sqrt(2)*sin(\x*180/pi + 45)});
\end{tikzpicture}

\end{document}
